\chapter{Удиртгал}

\section{Судалгааны үндэслэл ба орчин үеийн хэрэгцээ шаардлага}
Дэлхий даяар жижиглэн худалдааны (retail) салбар нь эдийн засгийн маш том хэсгийг бүрдүүлдэг бөгөөд хэрэглэгчдийн өдөр тутмын хэрэгцээт бараа бүтээгдэхүүнийг түгээх гол суваг болдог. Орчин үеийн жижиглэн худалдааны салбарт өрсөлдөөн улам бүр ширүүсч, үйлчлүүлэгчдэд хүргэх үйлчилгээний хурд, барааны олдоц, сонголт зэрэг нь тухайн дэлгүүрийн орлого, ашигт ажиллагаанд шууд нөлөөлөхүйц хүчин зүйлс болж байна. Энэхүү өрсөлдөөнт орчинд хамгийн чухал хэсгүүдийн нэг нь бараа материалын үр дүнтэй менежмент (Inventory Management), түүн дотроо лангуун дээрх барааны өрөлт, тооллого, татан авалтын хяналт юм.

Уламжлалт жижиглэн худалдааны менежментэд дэлгүүрийн ажилчид эсвэл бөөний төвийн борлуулалтын төлөөлөгчид (merchandiser) өөрсдийн биеэр лангууг шалгаж, ямар бараа дууссан, ямар барааны орлого нөхөгдөх шаардлагатай байгааг цаасан дээр болон гар утасны аппликейшн дээр гараар шивж бүртгэдэг байв. Энэхүү гар ажиллагаа шаардсан үйл явц нь маш их хугацаа зарцуулдаг төдийгүй, хүний хүчин зүйлээс шалтгаалсан алдаа гарах магадлал өндөр байдаг. Жишээлбэл, барааны үлдэгдлийг буруу тоолох, дууссан барааг хугацаанд нь мэдээлж чадахгүйгээс үүдэн борлуулалт алдах (Out-of-Stock loss), буруу бараа захиалах зэрэг асуудлууд байнга тохиолддог. Иймд бараа бүтээгдэхүүний өрөлтийн хяналтыг автоматжуулах, мэдээллийн үнэн зөв байдлыг сайжруулах хэрэгцээ улам бүр нэмэгдэж байна.

Сүүлийн арван жилд Хиймэл оюун ухаан (Artificial Intelligence - AI), тэр дундаа Компьютерын хараа (Computer Vision - CV) болон Гүнгэрвэл сургалт (Deep Learning) эрчимтэй хөгжихийн хэрээр дээрх асуудлуудыг технологийн хүчээр бүрэн шийдвэрлэх боломж бүрдлээ. Ухаалаг утасны камерын нягтрал сайжрахын зэрэгцээ серверийн тооцоолох хүчин чадал эрс нэмэгдсэн нь зөвхөн гэрэл зургийн мэдээлэлд тулгуурлан нарийн шинжилгээ хийж, дүгнэлт гаргах нөхцөлийг бий болгосон. Бид хамгийн сүүлийн үеийн алгоритмуудын нэг болох YOLO (You Only Look Once) архитектурыг ашигласнаар дэлгүүрийн лангууны зурган дээр байрлах бараануудыг богино хугацаанд, өндөр нарийвчлалтайгаар таньж, бүртгэх боломжтойг энэхүү судалгааны ажлаараа харуулахыг зорилоо. 

\section{Судалгааны тулгамдсан асуудал}
Бараа материалын менежментийг системчлэхэд дараах хэд хэдэн тулгамдсан асуудлууд оршиж байна:
\begin{enumerate}
    \item \textbf{Хүний нөөцийн зардал ба цаг хугацааны алдагдал:} Том хэмжээний супермаркетуудад олон мянган нэр төрлийн бараа байдаг бөгөөд эдгээрийг өдөр бүр тоолж шалгах нь асар их хөдөлмөр, хугацаа шаардсан ажил байдаг. Энэ нь ажилчдын үндсэн ажлын цагаас ихээхэн хумсалдаг.
    \item \textbf{Хүний хүчин зүйлээс үүдэлтэй алдаа (Human Error):} Олон цагаар нэгэн хэвийн тооллого хийх үед ажилтан ядарч, барааг дутуу тоолох эсвэл өөр төстэй бараатай андуурах тохиолдол элбэг байдаг.
    \item \textbf{Бодит хугацааны аналитик мэдээллийн дутагдал:} Уламжлалт хэлбэрээр хийсэн аудитын мэдээлэл төв системд орох хүртэл хэдэн цаг эсвэл өдрөөр саатдаг тул дээд удирдлага бодит хугацаанд шуурхай шийдвэр гаргах боломжгүйгээр хязгаарлагддаг.
    \item \textbf{Үнэтэй тоног төхөөрөмжийн хамаарал:} Гадны өндөр хөгжилтэй орнуудад нэвтэрч буй зарим автоматжуулалтын системүүд (жишээлбэл лангуун дээрх жин мэдрэгч, IoT камерууд г.м) нь суурилуулалтын анхны өртөг маш өндөр байдгаас хөгжиж буй орнуудын жижиглэн худалдааны зах зээлд хэрэглэхэд тохиромжгүй байдаг.
\end{enumerate}

Дээрх асуудлуудыг шийдвэрлэхийн тулд ямар нэгэн үнэтэй тоног төхөөрөмж шаардахгүйгээр, зөвхөн борлуулалтын ажилтны байнга ашигладаг ухаалаг утасны камерт суурилсан хүчирхэг, хурдан, уян хатан "Аудитын шийдвэр гаргах автомат систем"-ийг боловсруулах сорилт бидний өмнө тулгарч байна.

\section{Судалгааны зорилго}
Энэхүү ажлын гол зорилго нь хөдөлгөөнт төхөөрөмжөөс (ухаалаг гар утаснаас) дарсан дэлгүүрийн лангууны гэрэл зурагт гүн сургалтын объект илрүүлэх технологиор анализ хийж, зах зээлд борлуулагдаж буй бараа бүтээгдэхүүнийг автоматаар таньж, улмаар тухайн лангуун дээрх барааны өрөлт зөв байгаа эсэх, тоо ширхэг нь үнэн зөв байгаа эсэхэд тооллогын ба аудитын шийдвэр гаргалтыг автоматжуулах төгсгөл хоорондын (End-to-End) цогц ухаалаг системийг зохион бүтээхэд оршино.

\section{Судалгааны зорилтууд}
Дээрх ерөнхий зорилгыг бүрэн дүүрэн биелүүлэхийн тулд дараах тодорхой зорилтуудыг дэвшүүлэн хэрэгжүүлэв:
\begin{itemize}
    \item \textbf{Онолын судалгаа ба шинжилгээ хийх:} Компьютерын харааны объектод суурилсан илрүүлэлтийн алгоритмуудын хөгжлийн чиг хандлага, тэдгээрийн архитектурыг гүнзгийрүүлэн судлах, жижиглэн худалдааны салбарт ашиглагдаж буй автоматжуулалтын системүүдийн давуулон сул талыг харьцуулан шинжлэх.
    \item \textbf{Өгөгдөл бүрдүүлэх ба загвар сургах:} Бодит орчинд лангууны барааны зургийн өгөгдлийн санг цуглуулж, объект бүрийг тэмдэглэх (annotation) замаар датасет бэлтгэх. Улмаар YOLOv8 (You Only Look Once хувилбар 8) архитектурыг ашиглан жижиглэн худалдааны барааг өндөр нарийвчлалтайгаар таних гүн сургалтын загвар бэлтгэж, сургалт (training) хийх.
    \item \textbf{Серверийн хөгжүүлэлт (Backend):} FastAPI хүрээг (framework) ашиглан зураг хүлээн авч, сургасан модель руугаа дамжуулан объектуудыг илрүүлэх, MongoDB өгөгдлийн сантай харьцаж илрүүлсэн үр дүнгээ хадгалах өндөр хурдны микро-сервис (REST API) хөгжүүлэх.
    \item \textbf{Автомат аудитын алгоритм боловсруулах:} Зургаас таньсан үр дүнг тухайн дэлгүүрт байвал зохих хүлээгдэж буй (expected) барааны мэдээлэлтэй харьцуулан автоматаар зөрүүг тооцоолж, аудитын төлөв (PASS, WARNING, FAIL) зэргийг математик загварчлалын дагуу гаргаж өгөх шийдвэр гаргах логик (Audit Engine) боловсруулах.
    \item \textbf{Гар утасны аппликейшн хөгжүүлэлт (Frontend Mobile):} Flutter технологийг ашиглан хэрэглэгч ашиглахад хялбар, iOS болон Android үйлдлийн системүүд дээр зэрэг ажиллах боломжтой, үр дүнгээ хурдан хугацаанд хүлээн авч интерактив байдлаар харах мобайл аппликейшн бүтээх.
\end{itemize}

\section{Судалгааны обьект ба хамрах хүрээ}
Энэхүү судалгааны обьект нь дэлгүүрийн лангуун дээрх амьд гэрэл зураг ба түүн дээр хийгдэх компьютерын харааны илрүүлэлтийн үйл явц юм. Хамрах хүрээний хувьд:
\begin{itemize}
    \item Бүх төрлийн барааг хамрахаас илүүтэйгээр туршилтын зорилгоор дотоод гадаад худалдааны өргөн хэрэглээний ундаа, ус, жүүс зэрэг багц бараанууд (FMCG - Fast Moving Consumer Goods)-ийг онцолж датасетаа бүрдүүлсэн.
    \item Хөгжүүлсэн систем нь интернет сүлжээнд холбогдсон гар утсаар дамжуулан क्लाуд-д (Cloud) суурилсан сервертэй холбогдон ажиллахад зориулагдсан. Офлайн орчинд гар утас дээр шууд AI тооцоолол хийх (Edge computing) асуудал энэ удаагийн хамрах хүрээнээс гадуур болно.
\end{itemize}

\section{Судалгааны шинжлэх ухааны болон практикийн ач холбогдол}
\textbf{Шинжлэх ухааны ач холбогдол:}
Компьютерын харааны гүн сургалтын загварууд (ялангуяа YOLOv8) нь маш олон төрлийн саадтай (гэрэлтүүлэг, хагас халхлагдсан байдал г.м) бодит орчинд барааг таних нарийвчлалыг хэрхэн хадгалж байгааг онолын болон туршилтын үүднээс баталгаажуулав. Мөн машин сургалтын загварыг микросервисийн орчинд үр дүнтэй байршуулж, бусад системүүдтэй хэрхэн өндөр хурдтай уялдан ажиллах архитектурын шийдлийг гаргасан.

\textbf{Практикийн ач холбогдол:}
Санхүү болон цаг хугацааны хувьд асар их хэмнэлтийг бизнесийн байгууллагуудад авчирна. Дэлгүүрийн ажилтан эсвэл түгээгч нь зөвхөн цөөхөн хэдэн зураг дарахад л тооллогыг маш хурдан хийж, удирдлагууд дэлгүүрүүдийнхээ бодит нөхцөл байдлыг цаг алдалгүйгээр гар утас болон төвлөрсөн системээс хянах бүрэн боломжийг олгож байгаагаараа практикт шууд нэвтрэх өндөр нөөц бололцоотой юм.

\section{Ажлын бүтэц}
Энэхүү дипломын ажил нь нийтдээ 6 бүлгээс бүрдэх бөгөөд бүлэг тус бүрийн агуулга дараах байдалтай байна:
\begin{itemize}
    \item \textbf{Нэгдүгээр бүлэг (Удиртгал):} Судалгааны үндэслэл, ач холбогдол, зорилго, зорилтууд болон судалгааны объектын талаарх дэлгэрэнгүй танилцуулга.
    \item \textbf{Хоёрдугаар бүлэг (Онолын судалгаа ба ижил төстэй системүүд):} Компьютерын хараа болон объект илрүүлэлтийн алгоритмуудын үүсэл хөгжил, тэдгээрийн ялгаа, мөн дэлхийн жижиглэн худалдааны зах зээл дэх ижил төстэй ухаалаг системүүдийн судалгаа, харьцуулалт.
    \item \textbf{Гуравдугаар бүлэг (Системийн зохиомж ба арга зүй):} Төслийн ерөнхий микросервис архитектур, YOLOv8 загварын дотоод бүтэц, аудитын логик шийдвэр гаргалтын математик загварчлал зэргийг онолын үүднээс тайлбарлана.
    \item \textbf{Дөрөвдүгээр бүлэг (Системийн хэрэгжүүлэлт):} Өгөгдлийн сангийн бүрдүүлэлт, загварын сургалт хийсэн туршлага, FastAPI Backend болон Flutter Frontend аппликейшний хөгжүүлэлтийн техникийн нарийвчилсан мэдээллийг хүргэнэ.
    \item \textbf{Тавдугаар бүлэг (Үр дүн ба үнэлгээ):} Сургасан ML загварын нарийвчлалын үнэлгээ (mAP, Precision, Recall), системийн боловсруулах хурд, алдааны анализ, бодит амьдрал дээрх туршилтын үр дүн.
    \item \textbf{Зургаадугаар бүлэг (Дүгнэлт ба цаашдын зорилт):} Судалгааны ажлаас гарсан гол дүгнэлтүүд, системийн хязгаарлагдмал байдал болон цаашид хөгжүүлэх боломжуудын талаарх санал.
\end{itemize}
